\chapter{Codici MATLAB: Controllo Ottimo LQR}

In questa appendice vengono riportati i codici MATLAB utilizzati. Il documento è strutturato per essere facilmente navigabile: partendo dallo script principale, è possibile esplorare tutte le funzioni richiamate.

\section*{Script Principale LQR}
\subsection*{File: \texttt{LQRAnalysis.m}}
\label{sec:LQRAnalysis}

Script MATLAB principale per l'analisi di stabilità, tuning dei pesi tramite regola di Bryson e progettazione del controllo LQR. In questo script vengono richiamate le seguenti funzioni, definite nelle sezioni successive:
\begin{itemize}
    \item \hyperref[sec:plot_lyapunov_lqr]{\texttt{plot\_lyapunov\_lqr.m}}
    \item \hyperref[sec:funzioni_lqr_progetta_LQR_continuo]{\texttt{funzioni\_lqr/progetta\_LQR\_continuo.m}}
    \item \hyperref[sec:funzioni_lqr_progetta_LQR_discreto]{\texttt{funzioni\_lqr/progetta\_LQR\_discreto.m}}
    \item \hyperref[sec:funzioni_lqr_plot_lyapunov]{\texttt{funzioni\_lqr/plot\_lyapunov.m}}
    \item \hyperref[sec:funzioni_lqr_plot_lyapunov_discrete]{\texttt{funzioni\_lqr/plot\_lyapunov\_discrete.m}}
\end{itemize}

\begin{lstlisting}[language=Matlab, caption={Script Principale LQR}, label={lst:LQRAnalysis}]
% =========================================================================
% Tesi Triennale - LQR Longitudinale F-16
% Ing. Leggeri Leonardo
% =========================================================================
close all
% Inizializzazione Path
startup_project();

%% 1. IMPOSTAZIONE DEL SISTEMA (Solo veri attuatori)
% Estraiamo le dimensioni dalle matrici
nx = size(A_long, 1);
nu_long = size(B_long, 2);

B_ctrl = B_long(:, 1:3); 
D_ctrl = D_long(:, 1:3);
nu = size(B_ctrl, 2); 

%% 2. SINTESI DEL CONTROLLORE LQR 
% =========================================================================
% NOTA SUI PESI Q e R:
% Se vuoi modificare l'aggressività del controllore (pesi Q e R),
% apri il file "funzioni_lqr/progetta_LQR_continuo.m" e modificali 
% =========================================================================
disp('--- Sintesi LQR Continuo ---');
[K, P, Q, R, A_cl] = progetta_LQR_continuo(A_long, B_ctrl);

disp('Guadagno K:');
disp(K);
disp('Autovalori a ciclo chiuso (devono avere parte reale negativa):'); disp(eig(A_cl));

% Creiamo il sistema State-Space a ciclo chiuso (Continuo)
sys_cl = ss(A_cl, zeros(nx, nu), eye(nx), zeros(nx, nu));

% Vettore tempo per la simulazione
t = 0:0.05:10;

% Condizione iniziale estrema (Looping/Candela)
x0 = [0; 0.3; 0; 20]; % theta, q, U, W

% Simulazione della risposta libera (initial) del sistema
[y_sim, t_sim, x_sim] = initial(sys_cl, x0, t);

% Traspongo x_sim per avere gli stati sulle righe (come nel tuo script precedente)
xf = x_sim'; 

% Calcolo dei comandi u generati dall'LQR in ogni istante (u = -K*x)
uf = -K * xf; 

%% 4. GRAFICI (Senza vincoli, con nuovi colori)
rad2deg = 180 / pi; % Fattore di conversione utile per la leggibilità

% Definizione del colore Giallo leggibile per sfondo bianco
colore_giallo = [0.9290, 0.6940, 0.1250]; 

% -------------------------------------------------------------------------
% FIGURA 1: DINAMICA DEGLI STATI (Variabili Separate)
% -------------------------------------------------------------------------
figure('Name', 'Controllo LQR Puro: Stati a Ciclo Chiuso', 'Color', 'w', 'Position', [100, 100, 800, 800])

% 1. Angolo di Beccheggio (theta)
subplot(4, 1, 1)
plot(t, xf(1,:) * rad2deg, 'b', 'LineWidth', 2); hold on;
yline(0, 'k--', 'LineWidth', 1);
ylabel('Ampiezza [deg]', 'Interpreter', 'latex')
title('$\theta$ (Pitch Angle)', 'Interpreter', 'latex')
grid on;

% 2. Rateo di Beccheggio (q)
subplot(4, 1, 2)
plot(t, xf(2,:) * rad2deg, 'r', 'LineWidth', 2); hold on;
yline(0, 'k--', 'LineWidth', 1);
ylabel('Ampiezza [deg/s]', 'Interpreter', 'latex')
title('$q$ (Pitch Rate)', 'Interpreter', 'latex')
grid on;

% 3. Velocità asse X (U)
subplot(4, 1, 3)
plot(t, xf(3,:), 'Color', colore_giallo, 'LineWidth', 2); hold on;
yline(0, 'k--', 'LineWidth', 1);
ylabel('Ampiezza [ft/s]', 'Interpreter', 'latex')
title('$U$ (Velocit\`a Orizzontale)', 'Interpreter', 'latex')
grid on;

% 4. Velocità asse Z (W)
subplot(4, 1, 4)
plot(t, xf(4,:), 'g', 'LineWidth', 2); hold on;
yline(0, 'k--', 'LineWidth', 1);
xlabel('Tempo [s]', 'Interpreter', 'latex') 
ylabel('Ampiezza [ft/s]', 'Interpreter', 'latex')
title('$W$ (Velocit\`a Verticale $\approx \alpha$)', 'Interpreter', 'latex')
grid on;

% -------------------------------------------------------------------------
% FIGURA 2: SFORZO DEGLI ATTUATORI (Senza linee di saturazione)
% -------------------------------------------------------------------------
figure('Name', 'Controllo LQR Puro: Sforzo Attuatori', 'Color', 'w', 'Position', [150, 150, 800, 800])

% 1. Thrust (Spinta del Motore)
subplot(3, 1, 1)
plot(t, uf(1,:), 'b', 'LineWidth', 2); hold on; % Blu
yline(0, 'k--', 'LineWidth', 1);
ylabel('Spinta [lbs]')
title('Azione di Controllo Libera: Manetta (Thrust)')
grid on;

% 2. Elevator (Equilibratore)
subplot(3, 1, 2)
plot(t, uf(2,:) * rad2deg, 'r', 'LineWidth', 2); hold on; % Rosso
yline(0, 'k--', 'LineWidth', 1);
ylabel('Deflessione [deg]')
title('Azione di Controllo Libera: Equilibratore')
grid on;

% 3. Leading Edge Flap (LEF)
subplot(3, 1, 3)
plot(t, uf(3,:) * rad2deg, 'g', 'LineWidth', 2); hold on; % Verde
yline(0, 'k--', 'LineWidth', 1);
xlabel('Tempo [s]') 
ylabel('Deflessione [deg]')
title('Azione di Controllo Libera: Leading Edge Flap')
grid on;

% =========================================================================
% VISUALIZZAZIONE GRAFICA DEI NUOVI POLI (Ciclo Chiuso)
% =========================================================================
sys_cl_map = ss(A_cl, B_ctrl, C_long, D_ctrl); 

figure('Name', 'Mappa Poli-Zeri a Ciclo Chiuso (LQR Continuo)', 'Color', 'w', 'Position', [200, 200, 600, 500]);
pzmap(sys_cl_map);
grid on;

title('Verifica Stabilità LQR: Posizione dei Nuovi Poli');
xlabel('Asse Reale (Velocità di convergenza/divergenza)');
ylabel('Asse Immaginario (Frequenza di oscillazione)');

%% 5. RITRATTO DI FASE E FUNZIONE DI LYAPUNOV (LQR)
% Per visualizzare la stabilità, isoliamo la dinamica a corto periodo 
% variando W (velocità verticale) e q (pitch rate) e fissando U=0, theta=0.

% 5.1 Definizione della griglia spaziale
w_range = linspace(-25, 25, 50);   % ft/s
q_range_rad = linspace(-1.5, 1.5, 50); % rad/s
[W_grid, Q_rad] = meshgrid(w_range, q_range_rad);
Q_deg = Q_rad * rad2deg; % Conversione per asse grafico

% 5.2 Calcolo della funzione di Lyapunov V(x) = x^T * P * x sulla griglia
V_surf = zeros(size(W_grid));
for i = 1:size(W_grid, 1)
    for j = 1:size(W_grid, 2)
        % Fissiamo theta=0 e U=0
        stato_surf = [0; Q_rad(i,j); 0; W_grid(i,j)];
        V_surf(i,j) = stato_surf' * P * stato_surf;
    end
end

% 5.3 Dinamica a ciclo chiuso per ode45
dxdt = @(t, x) A_cl * x;

%% 5. RITRATTO DI FASE E FUNZIONE DI LYAPUNOV 3D (LQR)
% Usa la nuova funzione dedicata per generare il ritratto di fase 2D e
% la funzione di Lyapunov in 3D
plot_lyapunov_lqr(P, A_cl);
\end{lstlisting}

\section*{Analisi Funzionale di Lyapunov}
\subsection*{File: \texttt{plot\_lyapunov\_lqr.m}}
\label{sec:plot_lyapunov_lqr}

Script principale per il plotting, l'analisi e lo studio dell'ellissoide associato al funzionale di Lyapunov del sistema LQR.

\begin{lstlisting}[language=Matlab, caption={Analisi Funzionale di Lyapunov}, label={lst:plot_lyapunov_lqr}]
% =========================================================================
% Copyright (c) 2026 Ing. Leggeri Leonardo
% Tutti i diritti riservati.
%
% ATTENZIONE: Questo software e il relativo codice sorgente sono di proprietà 
% esclusiva dell'Ing. Leggeri Leonardo. È severamente vietata la copia, 
% la distribuzione, la modifica o la vendita a terzi senza l'esplicito 
% consenso scritto dell'autore. La distribuzione o la vendita non autorizzata 
% costituisce reato ed è perseguibile penalmente secondo le leggi vigenti.
% =========================================================================

function plot_lyapunov_lqr(P, A_cl)
    % =========================================================================
    % PLOT_LYAPUNOV - Ritratto di Fase e Funzione di Lyapunov 3D
    % Riceve in ingresso la matrice P e la matrice a ciclo chiuso A_cl
    % =========================================================================
    
    set(0,'DefaultLineLineWidth',1.5);
    set(0,'DefaultAxesFontSize',14);
    set(0,'DefaulttextInterpreter','latex');
    
    %% 1. Traiettorie (Solo dinamiche veloci)
    dxdt = @(t,x) A_cl * x;
    
    % Matrice delle condizioni iniziali multiple (8 scenari di partenza)
    % Ordine degli stati: [theta; q; u; w]
    x0_mult = [
         0,    0,    0,    0,    0,    0,    0,    0;
         1.2, -1.2,  0.5, -0.5,  0.8, -0.8,  0.3, -0.3;
         0,    0,    0,    0,    0,    0,    0,    0;
        20,  -20,   15,  -15, -10,   10,   22,  -22
    ]; 
          
    figure('Name', 'Ritratto di Fase LQR', 'Color', 'w', 'Position', [50 100 800 600]);
    hold on; grid on;
    
    % Generiamo automaticamente una palette di colori per le N traiettorie
    num_simulazioni = size(x0_mult, 2);
    colori = lines(num_simulazioni); 
    
    for i = 1:num_simulazioni
        % Simuliamo solo per 3 secondi per il corto periodo
        [~, x] = ode45(dxdt, [0 3], x0_mult(:, i)); 
        
        % Salviamo gli handle solo della prima iterazione per una legenda pulita
        if i == 1
            h_traj = plot(x(:,4), rad2deg(x(:,2)), 'Color', colori(i,:));
            h_start = plot(x(1,4), rad2deg(x(1,2)), '*', 'Color', colori(i,:), 'MarkerSize', 8); 
        else
            plot(x(:,4), rad2deg(x(:,2)), 'Color', colori(i,:));
            plot(x(1,4), rad2deg(x(1,2)), '*', 'Color', colori(i,:), 'MarkerSize', 8); 
        end
    end
    
    %% 2. Campo Vettoriale e Curve di Livello
   
    % Riduciamo la griglia per "abbracciare" esattamente le tue condizioni iniziali
    w_lim = 45;  
    q_lim = 75;  
    [W_grid, Q_deg] = meshgrid(linspace(-w_lim, w_lim, 80), linspace(-q_lim, q_lim, 80));
   
    Q_rad = deg2rad(Q_deg);
    
    W_dot = zeros(size(W_grid));
    Q_dot = zeros(size(Q_rad));
    V_contour = zeros(size(W_grid));
    
    for i = 1:numel(W_grid)
        % Slicing: imponiamo theta=0 e u=0
        stato = [0; Q_rad(i); 0; W_grid(i)];
        derivata = A_cl * stato;
        
        Q_dot(i) = derivata(2); 
        W_dot(i) = derivata(4); 
        V_contour(i) = stato' * P * stato;
    end
    
    % Campo vettoriale
    L = sqrt(W_dot.^2 + Q_dot.^2) + 1e-6;
    h_quiv = quiver(W_grid, Q_deg, W_dot./L, rad2deg(Q_dot)./L, 0.45, 'Color', [0.7 0.7 0.7]);
    
    % Curve di Livello (Spaziatura Logaritmica)
    val_min = min(V_contour(:));
    val_max = max(V_contour(:));
    livelli_log = logspace(log10(val_min + 1e-3), log10(val_max), 18);
    
    [~, h_cont] = contour(W_grid, Q_deg, V_contour, livelli_log, 'LineWidth', 1.5);
    
    xlabel('Vertical Velocity $w$ [ft/s]') 
    ylabel('Pitch Rate $q$ [deg/s]') 
    title('Ritratto di fase')
    
    xlim([-w_lim, w_lim]); 
    ylim([-q_lim, q_lim]);
    
    % Legenda dinamica e pulita (solo 4 voci invece di 18)
    legend([h_traj, h_start, h_quiv, h_cont], ...
           {'Traiettorie di Stato', 'Punti di Partenza', 'Direzione Flusso', 'Contour $V(x) = x^T P x$'}, ...
           'Interpreter','latex', 'Location', 'bestoutside');

    %% 3. Visualizzazione 3D della Funzione di Lyapunov LQR
    figure('Name', 'Funzione di Lyapunov 3D - LQR', 'Color', 'w');
    hold on; grid on;
    
    % Superficie 3D
    V_surf = zeros(size(W_grid));
    for i = 1:size(W_grid, 1)
        for j = 1:size(W_grid, 2)
            stato_surf = [0; Q_rad(i,j); 0; W_grid(i,j)];
            V_surf(i,j) = stato_surf' * P * stato_surf;
        end
    end
    
    h_surf = surf(W_grid, Q_deg, V_surf, 'EdgeColor', 'none', 'FaceAlpha', 0.7);
    colormap jet;
    colorbar;
    
    % --- AGGIUNTA FONDAMENTALE: Inizializziamo max_v_traj ---
    max_v_traj = 0;
    
    % Proiezione traiettorie 3D
    for i = 1:num_simulazioni
        % Integrazione più fitta per fluidità visiva
        [~, x] = ode45(dxdt, linspace(0, 3, 300), x0_mult(:, i));
        
        V_traiettoria = sum((x * P) .* x, 2); 
        
        % --- Calcoliamo il picco di energia ---
        if max(V_traiettoria) > max_v_traj
            max_v_traj = max(V_traiettoria);
        end
        
        % Usiamo i colori estratti dalla palette per mantenere coerenza col 2D
        if i == 1
            h_traj3 = plot3(x(:,4), rad2deg(x(:,2)), V_traiettoria, '-', 'Color', colori(i,:), 'LineWidth', 2.5);
            h_start3 = plot3(x(1,4), rad2deg(x(1,2)), V_traiettoria(1), 'o', 'MarkerFaceColor', colori(i,:), 'MarkerEdgeColor', 'k', 'MarkerSize', 6);
        else
            plot3(x(:,4), rad2deg(x(:,2)), V_traiettoria, '-', 'Color', colori(i,:), 'LineWidth', 2.5);
            plot3(x(1,4), rad2deg(x(1,2)), V_traiettoria(1), 'o', 'MarkerFaceColor', colori(i,:), 'MarkerEdgeColor', 'k', 'MarkerSize', 6);
        end
    end
    
    title('Funzione di Lyapunov $V(x) = x^T P x$ a Ciclo Chiuso', 'Interpreter', 'latex')
    xlabel('Velocit\`a Verticale $w$ [ft/s]', 'Interpreter', 'latex')
    ylabel('Pitch Rate $q$ [deg/s]', 'Interpreter', 'latex')
    zlabel('Energia $V(x)$', 'Interpreter', 'latex')
    
    % Taglio dinamico dell'asse Z per non tagliare il paraboloide
    zmax_plot = max(max(V_surf(:)), max_v_traj);
    zlim([0, zmax_plot * 1.1]); 
    xlim([-w_lim, w_lim]);
    ylim([-q_lim, q_lim]);
    view(-45, 30);
    
    % Legenda 3D
    legend([h_surf, h_traj3, h_start3], ...
           {'Funzione di Lyapunov $V(x)$', 'Traiettoria di Stato', 'Condizione Iniziale'}, ...
           'Interpreter', 'latex', 'Location', 'bestoutside');
end
\end{lstlisting}

\section*{Progetto LQR (Tempo Continuo)}
\subsection*{File: \texttt{funzioni\_lqr/progetta\_LQR\_continuo.m}}
\label{sec:funzioni_lqr_progetta_LQR_continuo}

Funzione per la risoluzione dell'equazione algebrica di Riccati (ARE) e il calcolo del guadagno ottimo LQR nel dominio del tempo continuo.

\begin{lstlisting}[language=Matlab, caption={Progetto LQR (Tempo Continuo)}, label={lst:funzioni_lqr_progetta_LQR_continuo}]
% =========================================================================
% Copyright (c) 2026 Ing. Leggeri Leonardo
% Tutti i diritti riservati.
%
% ATTENZIONE: Questo software e il relativo codice sorgente sono di proprietà 
% esclusiva dell'Ing. Leggeri Leonardo. È severamente vietata la copia, 
% la distribuzione, la modifica o la vendita a terzi senza l'esplicito 
% consenso scritto dell'autore. La distribuzione o la vendita non autorizzata 
% costituisce reato ed è perseguibile penalmente secondo le leggi vigenti.
% =========================================================================

function [K, P, Q, R, A_cl] = progetta_LQR_continuo(A, B)
    % =====================================================================
    % Calcolo pesi LQR (Regola di Bryson) coerenti con il modello F-16
    % Ordine stati: [theta (rad); q (rad/s); U (ft/s); W (ft/s)]
    % Ordine input: [Thrust (lbs); Elevator (rad); LEF (rad)]
    % =====================================================================
    
    % Massimi scostamenti tollerabili per gli stati (Bryson's rule)
    max_theta = deg2rad(15); % 15 gradi max tollerati per il pitch
    max_q     = deg2rad(30); % 30 deg/s max tollerati per il pitch rate
    max_U     = 30;          % 30 ft/s max per la velocità X
    max_W     = 20;          % 20 ft/s max per la velocità Z (incide su alpha)
    
    Q = diag([1/max_theta^2, 1/max_q^2, 1/max_U^2, 1/max_W^2]);
    
    % Massimi scostamenti tollerabili per gli attuatori
    max_thrust = 5000;       % 5000 lbs variazione manetta
    max_elev   = deg2rad(25);% 25 gradi max deflessione equilibratore
    max_lef    = deg2rad(25);% 25 gradi max deflessione LEF
    
    R = diag([1/max_thrust^2, 1/max_elev^2, 1/max_lef^2]);
    
    % Fattori di tuning per dare più aggressività globale (se necessario)
    rho_q = 1; 
    rho_r = 100; 
    
    Q = rho_q * Q;
    R = rho_r * R;
    
    % Sintesi LQR
    [K, P, ~] = lqr(A, B, Q, R);
    A_cl = A - B*K;
end
\end{lstlisting}

\section*{Progetto LQR}
\subsection*{File: \texttt{funzioni\_lqr/progetta\_LQR\_discreto.m}}
\label{sec:funzioni_lqr_progetta_LQR_discreto}

Funzione per il calcolo del guadagno ottimo LQR nel dominio del tempo discreto.

\begin{lstlisting}[language=Matlab, caption={Progetto LQR (Tempo Discreto)}, label={lst:funzioni_lqr_progetta_LQR_discreto}]
% =========================================================================
% Copyright (c) 2026 Ing. Leggeri Leonardo
% Tutti i diritti riservati.
%
% ATTENZIONE: Questo software e il relativo codice sorgente sono di proprietà 
% esclusiva dell'Ing. Leggeri Leonardo. È severamente vietata la copia, 
% la distribuzione, la modifica o la vendita a terzi senza l'esplicito 
% consenso scritto dell'autore. La distribuzione o la vendita non autorizzata 
% costituisce reato ed è perseguibile penalmente secondo le leggi vigenti.
% =========================================================================

function [K, P, Q, R, A_cl] = progetta_LQR_discreto(A, B)
    % =====================================================================
    % Calcolo pesi LQR (Regola di Bryson) coerenti con il modello F-16
    % Ordine stati: [theta (rad); q (rad/s); U (m/s); W (m/s)]
    % Ordine input: [Thrust (N); Elevator (rad); LEF (rad)]
    % =====================================================================
    
    ft2m = 0.3048;
    lbf2N = 4.44822;
    
    % Massimi scostamenti tollerabili per gli stati (Bryson's rule)
    max_theta = 15 * pi/180; % 15 gradi max tollerati per il pitch
    max_q     = 30 * pi/180; % 30 deg/s max tollerati per il pitch rate
    max_U     = 30 * ft2m;   % 30 ft/s convertiti in m/s max per la velocità X
    max_W     = 20 * ft2m;   % 20 ft/s convertiti in m/s max per la velocità Z
    
    Q = diag([1/max_theta^2, 1/max_q^2, 1/max_U^2, 1/max_W^2]);
    
    % Massimi scostamenti tollerabili per gli attuatori
    max_thrust = 5000 * lbf2N;   % 5000 lbs convertite in Newton
    max_elev   = 25 * pi/180;    % 25 gradi max deflessione equilibratore
    max_lef    = 25 * pi/180;    % 25 gradi max deflessione LEF
    
    R = diag([1/max_thrust^2, 1/max_elev^2, 1/max_lef^2]);
    
    % Fattori di tuning per dare più aggressività globale (se necessario)
    rho_q = 20;  % Bilanciamento: molto reattivo ma senza rendere il problema Infeasible
    rho_r = 1;   % Mantiene uno sforzo di controllo ragionevole per non violare i vincoli con x_iniziale elevato 
    
    Q = rho_q * Q;
    R = rho_r * R;
    
    % Sintesi LQR
    fprintf('DEBUG progetta_LQR_discreto:\n');
    fprintf('Size of A: %dx%d\n', size(A,1), size(A,2));
    fprintf('Size of B: %dx%d\n', size(B,1), size(B,2));
    fprintf('Size of Q before dlqr: %dx%d\n', size(Q,1), size(Q,2));
    fprintf('Size of R before dlqr: %dx%d\n', size(R,1), size(R,2));
    
    % Assicuriamoci che Q e R siano delle dimensioni esatte
    Nx = size(A, 1);
    Nu = size(B, 2);
    
    if size(Q, 1) ~= Nx || size(Q, 2) ~= Nx
        disp('ATTENZIONE: Q non è della dimensione corretta! Ricostruzione sicura in corso...');
        Q_new = eye(Nx);
        n_min = min(Nx, size(Q,1));
        Q_new(1:n_min, 1:n_min) = Q(1:n_min, 1:n_min);
        Q = Q_new;
    end
    
    if size(R, 1) ~= Nu || size(R, 2) ~= Nu
        disp('ATTENZIONE: R non è della dimensione corretta! Ricostruzione sicura in corso...');
        R = eye(Nu);
    end
    
    [K, P, ~] = dlqr(A, B, Q, R);
    A_cl = A - B*K;
end
\end{lstlisting}

\section*{Plot Lyapunov (Tempo Continuo)}
\subsection*{File: \texttt{funzioni\_lqr/plot\_lyapunov.m}}
\label{sec:funzioni_lqr_plot_lyapunov}

Funzione per la visualizzazione grafica dell'ellissoide del funzionale di Lyapunov continuo.

\begin{lstlisting}[language=Matlab, caption={Plot Lyapunov (Tempo Continuo)}, label={lst:funzioni_lqr_plot_lyapunov}]
% =========================================================================
% Copyright (c) 2026 Ing. Leggeri Leonardo
% Tutti i diritti riservati.
%
% ATTENZIONE: Questo software e il relativo codice sorgente sono di proprietà 
% esclusiva dell'Ing. Leggeri Leonardo. È severamente vietata la copia, 
% la distribuzione, la modifica o la vendita a terzi senza l'esplicito 
% consenso scritto dell'autore. La distribuzione o la vendita non autorizzata 
% costituisce reato ed è perseguibile penalmente secondo le leggi vigenti.
% =========================================================================

function plot_lyapunov(P, A_cl)
    % =========================================================================
    % PLOT_LYAPUNOV - Ritratto di Fase e Funzione di Lyapunov 3D
    % Riceve in ingresso la matrice P e la matrice a ciclo chiuso A_cl
    % =========================================================================
    
    set(0,'DefaultLineLineWidth',1.5);
    set(0,'DefaultAxesFontSize',14);
    set(0,'DefaulttextInterpreter','latex');
    
    %% 1. Traiettorie (Solo dinamiche veloci)
    dxdt = @(t,x) A_cl * x;
    
    % Condizioni iniziali 
    % Matrice delle condizioni iniziali multiple (8 scenari di partenza)
% Ordine degli stati: [theta; q; u; w]
x0_mult = [
     0,    0,    0,    0,    0,    0,    0,    0;
     1.2, -1.2,  0.5, -0.5,  0.8, -0.8,  0.3, -0.3;
     0,    0,    0,    0,    0,    0,    0,    0;
    20,  -20,   15,  -15, -10,   10,   22,  -22
]; 
          
    figure('Name', 'Corto Periodo: w vs q', 'Color', 'w', 'Position', [50 100 800 600]);
    hold on; grid on;
    colori = {'r', 'g', 'b'}; 
    
    for i = 1:size(x0, 2)
        % Simuliamo solo per 3 secondi per il corto periodo
        [t, x] = ode45(dxdt, [0 3], x0(:, i)); 
        % Plot: Asse X = w (ft/s), Asse Y = q (deg/s)
        plot(x(:,4), rad2deg(x(:,2)), 'Color', colori{i});
        plot(x(1,4), rad2deg(x(1,2)), '*', 'Color', colori{i}, 'MarkerSize', 8); 
    end

    %% 2. Campo Vettoriale e Curve di Livello
    w_lim = 150;
    q_lim = 150;
    
    [W_grid, Q_deg] = meshgrid(linspace(-w_lim, w_lim, 35), linspace(-q_lim, q_lim, 35));
    Q_rad = deg2rad(Q_deg);
    
    W_dot = zeros(size(W_grid));
    Q_dot = zeros(size(Q_rad));
    V_contour = zeros(size(W_grid));
    
    for i = 1:numel(W_grid)
        % Slicing: imponiamo theta=0 e u=0
        stato = [0; Q_rad(i); 0; W_grid(i)];
        derivata = A_cl * stato;
        
        Q_dot(i) = derivata(2); 
        W_dot(i) = derivata(4); 
        V_contour(i) = stato' * P * stato;
    end
    
    % Campo vettoriale
    L = sqrt(W_dot.^2 + Q_dot.^2) + 1e-6;
    quiver(W_grid, Q_deg, W_dot./L, rad2deg(Q_dot)./L, 0.45, 'Color', [0.7 0.7 0.7]);
    
    % Curve di Livello (Spaziatura Logaritmica)
    val_min = min(V_contour(:));
    val_max = max(V_contour(:));
    livelli_log = logspace(log10(val_min + 1e-3), log10(val_max), 18);
    
    contour(W_grid, Q_deg, V_contour, livelli_log, 'LineWidth', 1.5);
    xlabel('Vertical Velocity $w$ [ft/s]') 
    ylabel('Pitch Rate $q$ [deg/s]') 
    title('Moto di Corto Periodo (Dinamica Veloce)')
    
    xlim([-w_lim, w_lim]); 
    ylim([-q_lim, q_lim]);
    legend({'Traiettoria 1', 'Start 1', 'Traiettoria 2', 'Start 2', 'Traiettoria 3', 'Start 3', ...
            'Direzione Flusso', 'Contour $V(x) = x^T P x$'}, 'Interpreter','latex', 'Location', 'bestoutside');

    %% 3. Visualizzazione 3D della Funzione di Lyapunov LQR
    figure('Name', 'Funzione di Lyapunov 3D - LQR', 'Color', 'w');
    hold on; grid on;
    
    % Superficie 3D
    V_surf = zeros(size(W_grid));
    for i = 1:size(W_grid, 1)
        for j = 1:size(W_grid, 2)
            stato_surf = [0; Q_rad(i,j); 0; W_grid(i,j)];
            V_surf(i,j) = stato_surf' * P * stato_surf;
        end
    end
    
    surf(W_grid, Q_deg, V_surf, 'EdgeColor', 'none', 'FaceAlpha', 0.7);
    colormap jet;
    colorbar;
    
    % Proiezione traiettorie 3D
    for i = 1:size(x0, 2)
        % Integrazione più fitta per fluidità visiva
        [t, x] = ode45(dxdt, linspace(0, 3, 300), x0(:, i));
        
        V_traiettoria = sum((x * P) .* x, 2); 
        
        plot3(x(:,4), rad2deg(x(:,2)), V_traiettoria, 'k-', 'LineWidth', 2);
        plot3(x(1,4), rad2deg(x(1,2)), V_traiettoria(1), 'ro', 'MarkerFaceColor', 'r', 'MarkerSize', 6);
    end
    
    title('Funzione di Lyapunov $V(x) = x^T P x$ (LQR)', 'Interpreter', 'latex')
    xlabel('$w$ [ft/s]', 'Interpreter', 'latex')
    ylabel('$q$ [deg/s]', 'Interpreter', 'latex')
    zlabel('$V(x)$', 'Interpreter', 'latex')
    view(-45, 30);
end
\end{lstlisting}

\section*{Plot Lyapunov (Tempo Discreto)}
\subsection*{File: \texttt{funzioni\_lqr/plot\_lyapunov\_discrete.m}}
\label{sec:funzioni_lqr_plot_lyapunov_discrete}

Funzione per la visualizzazione grafica dell'ellissoide del funzionale di Lyapunov discreto.

\begin{lstlisting}[language=Matlab, caption={Plot Lyapunov (Tempo Discreto)}, label={lst:funzioni_lqr_plot_lyapunov_discrete}]
% =========================================================================
% Copyright (c) 2026 Ing. Leggeri Leonardo
% Tutti i diritti riservati.
%
% ATTENZIONE: Questo software e il relativo codice sorgente sono di proprietà 
% esclusiva dell'Ing. Leggeri Leonardo. È severamente vietata la copia, 
% la distribuzione, la modifica o la vendita a terzi senza l'esplicito 
% consenso scritto dell'autore. La distribuzione o la vendita non autorizzata 
% costituisce reato ed è perseguibile penalmente secondo le leggi vigenti.
% =========================================================================

function plot_lyapunov_discrete(P_ds, A_cl_ds, Ts)
    % =========================================================================
    % PLOT_LYAPUNOV_DISCRETE - Funzione di Lyapunov 3D per sistemi a tempo discreto
    % Grafico migliorato: Traiettoria interpolata fluida + Campionamenti discreti
    % =========================================================================
    
    if nargin < 3
        Ts = 0.05; 
    end
    
    rad2deg = 180 / pi;
    
    % --- 1. DEFINIZIONE LIMITI VISIVI ---
    w_lim = 50;  % ft/s
    q_lim = 50;  % deg/s
    
    [W_grid, Q_deg] = meshgrid(linspace(-w_lim, w_lim, 80), linspace(-q_lim, q_lim, 80));
    Q_rad = Q_deg / rad2deg; 
    
    % --- 2. CALCOLO SUPERFICIE DI LYAPUNOV ---
    V_surf = zeros(size(W_grid));
    for i = 1:size(W_grid, 1)
        for j = 1:size(W_grid, 2)
            stato_surf = [0; Q_rad(i,j); 0; W_grid(i,j)];
            V_surf(i,j) = stato_surf' * P_ds * stato_surf;
        end
    end
    
    % --- 3. DINAMICA DISCRETA E CONDIZIONI INIZIALI ---
    N_steps = round(10 / Ts); 
    
    x0 = [0,  0.5, 0,  30;   
          0, -0.6, 0, -40;   
          0,  0.8, 0, -20]';

    %% --- FIGURA 3D ---
    figure('Name', sprintf('Lyapunov 3D Discreto (Ts = %g s)', Ts), 'Color', 'w', 'Position', [150 150 850 650]);
    hold on; grid on;
    
    % 1. Disegno Superficie
    h_surf = surf(W_grid, Q_deg, V_surf, 'EdgeColor', 'none', 'FaceAlpha', 0.65);
    colormap jet;
    cb = colorbar;
    ylabel(cb, 'Energia $V(x_k)$', 'Interpreter', 'latex', 'FontSize', 12);
    
    colori = {'r', 'g', 'b'};
    max_v_traj = 0; 
    
    % 2. Simulazione e Disegno Traiettorie
    for i = 1:size(x0, 2)
        x_traj = zeros(4, N_steps);
        x_traj(:, 1) = x0(:, i);
        
        for k = 1:N_steps-1
            x_traj(:, k+1) = A_cl_ds * x_traj(:, k);
        end
        
        V_traj = zeros(N_steps, 1);
        for k = 1:N_steps
            V_traj(k) = x_traj(:, k)' * P_ds * x_traj(:, k);
        end
        
        if max(V_traj) > max_v_traj
            max_v_traj = max(V_traj);
        end
        
        % --- MAGIA VISIVA: INTERPOLAZIONE FLUIDA DELLA TRAIETTORIA ---
        % Usiamo 'pchip' (interpolazione cubica che non crea finti rimbalzi) 
        % per generare 10 volte più punti e rendere la curva morbidissima
        t_discrete = 1:N_steps;
        t_fine = linspace(1, N_steps, N_steps * 10);
        
        W_smooth = pchip(t_discrete, x_traj(4,:), t_fine);
        Q_smooth = pchip(t_discrete, x_traj(2,:) * rad2deg, t_fine);
        V_smooth = pchip(t_discrete, V_traj, t_fine);
        
        % 1. Disegniamo prima la SCIA FLUIDA (linea continua)
        if i == 1
            h_line = plot3(W_smooth, Q_smooth, V_smooth, '-', 'Color', colori{i}, 'LineWidth', 2);
        else
            plot3(W_smooth, Q_smooth, V_smooth, '-', 'Color', colori{i}, 'LineWidth', 2);
        end
        
        % 2. Disegniamo sopra i PUNTI DISCRETI (solo pallini, senza linee)
        plot3(x_traj(4,:), x_traj(2,:) * rad2deg, V_traj, 'o', ...
              'MarkerEdgeColor', colori{i}, 'MarkerFaceColor', 'w', 'MarkerSize', 4);
              
        % Partenza (Start)
        if i == 1
            h_start = plot3(x_traj(4,1), x_traj(2,1) * rad2deg, V_traj(1), 'o', ...
                            'MarkerFaceColor', colori{i}, 'MarkerEdgeColor', 'k', 'MarkerSize', 8);
        else
            plot3(x_traj(4,1), x_traj(2,1) * rad2deg, V_traj(1), 'o', ...
                  'MarkerFaceColor', colori{i}, 'MarkerEdgeColor', 'k', 'MarkerSize', 8);
        end
        
        % Arrivo (End)
        plot3(x_traj(4,end), x_traj(2,end) * rad2deg, V_traj(end), 's', ...
              'MarkerFaceColor', colori{i}, 'MarkerEdgeColor', 'k', 'MarkerSize', 8);
    end
    
    % --- ORIGINE ---
    h_end = plot3(0, 0, 0, 'p', 'MarkerFaceColor', 'y', 'MarkerEdgeColor', 'k', 'MarkerSize', 16);
    
    % --- FORMATTAZIONE E LEGENDA ---
    title(sprintf('\\textbf{Funzione di Lyapunov $V(x_k)$ a Tempo Discreto ($T_s = %g$ s)}', Ts), 'Interpreter', 'latex', 'FontSize', 16);
    xlabel('Velocit\`a Verticale $w$ [ft/s]', 'Interpreter', 'latex', 'FontSize', 12);
    ylabel('Pitch Rate $q$ [deg/s]', 'Interpreter', 'latex', 'FontSize', 12);
    zlabel('Energia $V(x_k) = x_k^T P_d x_k$', 'Interpreter', 'latex', 'FontSize', 12);
    
    legend([h_surf, h_line, h_start, h_end], ...
           {'Superficie $V(x_k)$', 'Evoluzione di Stato (Interpolata)', 'Partenza ($k=0$)', 'Origine ($k \to \infty$)'}, ...
           'Interpreter', 'latex', 'FontSize', 12, 'Location', 'northeast');
    
    zlim([0, max_v_traj * 1.2]); 
    xlim([-w_lim, w_lim]);
    ylim([-q_lim, q_lim]);
    view(-35, 30);
end
\end{lstlisting}

